\documentclass[11pt,a4paper]{article}
\usepackage{iftex}
\ifPDFTeX
  \usepackage[T1]{fontenc}
  \usepackage[utf8]{inputenc}
  \usepackage{lmodern}
\else
  \usepackage{fontspec}
  \setmainfont[BoldFont=lmroman10-bold.otf,ItalicFont=lmroman10-italic.otf,BoldItalicFont=lmroman10-bolditalic.otf]{lmroman10-regular.otf}
  \setmonofont{lmmono10-regular.otf}
\fi
\usepackage[polish]{babel}
\usepackage[margin=2.5cm]{geometry}
\usepackage{microtype}
\usepackage{amsmath,amssymb}
\usepackage{booktabs}
\usepackage{hyperref}
\hypersetup{colorlinks=true,linkcolor=blue,urlcolor=blue}
\setlength{\parindent}{0pt}
\setlength{\parskip}{0.55em}
\newcommand{\doUzupelnienia}[1]{\textbf{Do uzupełnienia:} #1}
\title{Uczenie się ze wzmocnieniem\\Sprawozdanie z projektu}
\author{Autor: do uzupełnienia przez studenta}
\date{Wersja robocza: 21 października 2026}
\begin{document}
\maketitle
\textbf{Status dokumentu.} Wersja robocza rozwijana równolegle z projektem.
Temat „Górnik w niestabilnej kopalni” jest propozycją zgłaszaną do akceptacji
prowadzącego (etap E1). Formalizacja jest wstępna i może zostać skorygowana po
pierwszych testach środowiska. Nie przeprowadzono jeszcze eksperymentów.
\tableofcontents

\section{Opis problemu}
Górnik schodzi do kopalni, aby wydobyć złoża i wynieść je na powierzchnię.
Złoża położone głębiej są cenniejsze, ale każde wydobycie osłabia strop kopalni
i zwiększa ryzyko zawału. Lampa górnika ma ograniczony zapas oleju, a do
wzmacniania stropu służy niewielka liczba stempli. Zawał lub wyczerpanie oleju
oznacza utratę całego niesionego ładunku. Nagrodę agent otrzymuje dopiero wtedy,
gdy wyjdzie z kopalni.

Celem agenta jest maksymalizacja oczekiwanej wartości ładunku wyniesionego na
powierzchnię. W praktyce sprowadza się to do pytania, które agent musi sobie
zadawać w każdym kroku: \emph{kopać dalej czy wracać?} Projekt realizuje
podstawową ścieżkę środowiska gridowego (P00, sekcja 6.1) z energią (olej),
zasobami (stemple), zdarzeniami zależnymi od decyzji agenta oraz kilkoma mapami
o różnej trudności.

\section{Uzasadnienie zastosowania RL}
\begin{itemize}
  \item \textbf{Decyzje sekwencyjne.} Wynik epizodu zależy od całej sekwencji
  ruchów, wydobyć i stawiania stempli, a nie od pojedynczego wyboru.
  \item \textbf{Wpływ na przyszły stan.} Każde wydobycie trwale obniża stabilność
  stropu i zmienia mapę, a po wydobyciu określonej liczby złóż zawala się wybrany
  korytarz. Decyzja podjęta na początku epizodu wpływa na ryzyko i dostępne drogi
  w dalszej części.
  \item \textbf{Niepewność.} Zawał jest zdarzeniem losowym, którego
  prawdopodobieństwo rośnie wraz ze spadkiem stabilności i głębokością.
  Agent musi nauczyć się oceniać ryzyko, a nie tylko szukać najkrótszej drogi.
  \item \textbf{Opóźniona nagroda i konflikt celów.} Nagroda pojawia się dopiero
  przy wyjściu. Głębsze złoża zwiększają zysk, ale także ryzyko i czas powrotu,
  co wymaga planowania długoterminowego.
  \item \textbf{Dlaczego nie zwykłe przeszukiwanie.} Algorytmy takie jak BFS czy
  A* znajdują najkrótszą drogę w znanym, deterministycznym grafie. Tutaj
  optymalna strategia zależy od rozkładu prawdopodobieństwa zawału i od
  kompromisu między zyskiem a ryzykiem, który agent ma odkryć przez interakcję.
\end{itemize}

\section{Formalizacja RL/MDP}
\subsection{Agent i środowisko}
Agentem jest górnik. Środowiskiem jest kopalnia w postaci siatki
o szerokości $W$ i głębokości $H$ (wstępnie $W=8$, $H=10$). Wiersz $0$ to
powierzchnia z polem wyjścia. Pola siatki: lita skała (nieprzechodnia),
korytarz, złoże (węgiel, srebro, złoto, kryształ), wyjście oraz pole startowe.

\subsection{Stan}
Stan w kroku $t$ to krotka
\[
  s_t = (x,\ y,\ \sigma,\ o,\ m,\ n),
\]
gdzie $(x,y)$ to pozycja górnika, $\sigma \in \{0,\dots,5\}$ to globalna
stabilność stropu, $o \in \{0,1,2,3\}$ to zapas oleju w czterech przedziałach,
$m \in \{0,1\}^K$ to maska wydobytych złóż ($K \le 5$ złóż na mapie),
a $n \in \{0,1,2\}$ to liczba pozostałych stempli.

Wartość niesionego ładunku nie jest osobnym elementem stanu, ponieważ wynika
jednoznacznie z maski $m$ (ładunek to suma wartości wydobytych złóż). Zawalenie
korytarza zależy wyłącznie od liczby wydobytych złóż, czyli także od $m$.
Dzięki temu stan zachowuje własność Markowa, a jego rozmiar jest ograniczony do
$80 \cdot 6 \cdot 4 \cdot 2^5 \cdot 3 \approx 1{,}8 \cdot 10^5$ kombinacji,
z których osiągalna jest tylko część. Pozwala to na tablicową reprezentację
funkcji wartości.

\subsection{Akcje}
$\mathcal{A} = \{\text{góra}, \text{dół}, \text{lewo}, \text{prawo}, \text{podepnij}\}$.
\begin{itemize}
  \item Ruch na korytarz, wyjście lub pole po wydobytym złożu przesuwa górnika.
  \item Ruch na niewydobyte złoże oznacza jego wydobycie: górnik wchodzi na pole,
  ładunek rośnie o wartość złoża, a stabilność spada o 1 (nie poniżej 0).
  \item Ruch w litą skałę lub poza mapę jest akcją niedozwoloną: pozycja się nie
  zmienia, agent otrzymuje karę.
  \item \emph{Podepnij} zużywa jeden stempel i podnosi stabilność o 2
  (maksymalnie do 5). Bez stempli lub przy pełnej stabilności jest akcją
  niedozwoloną.
\end{itemize}
Świadomie zrezygnowano z akcji \emph{czekaj}: stabilność nie odbudowuje się
samoczynnie, więc czekanie nie przynosiłoby korzyści, a jedynie zużywało olej.

\subsection{Dynamika przejść}
W każdym kroku zapas oleju maleje o 1 (wartość początkowa 100). Po wykonaniu
akcji następuje losowanie zawału z prawdopodobieństwem
\[
  p_{\text{zawał}}(\sigma, y) = p_0 \,(5 - \sigma)\left(1 + \frac{y}{H}\right),
  \qquad p_0 = 0{,}02,
\]
gdzie $y$ to głębokość górnika. Przy pełnej stabilności zawał nie występuje,
a przy $\sigma = 0$ na dnie kopalni prawdopodobieństwo wynosi $0{,}2$ na krok.
Po wydobyciu $N=3$ złóż zawala się korytarz wskazany w opisie mapy, co może
odciąć najkrótszą drogę powrotu.

\subsection{Funkcja nagrody}
\begin{center}
\begin{tabular}{@{}ll@{}}
\toprule
Zdarzenie & Nagroda \\
\midrule
każdy krok & $-0{,}1$ \\
akcja niedozwolona & $-0{,}5$ (dodatkowo) \\
wyjście z kopalni & $+$wartość ładunku \\
zawał lub wyczerpanie oleju & $-20$, ładunek przepada \\
\bottomrule
\end{tabular}
\end{center}
Wartości złóż: węgiel 1, srebro 3, złoto 6, kryształ 10.
Koszt kroku zniechęca do zbędnego krążenia. Kara za niedozwoloną akcję
przyspiesza naukę ograniczeń mapy. Stała kara $-20$ sprawia, że porażka jest
dotkliwa także wtedy, gdy agent nie niesie ładunku. W eksperymentach zostanie
zbadany wpływ jej wartości (np. $-5$, $-20$, $-50$) na skłonność agenta do ryzyka.

\subsection{Warunki zakończenia epizodu}
Epizod kończy się, gdy górnik wyjdzie z kopalni, nastąpi zawał albo skończy się
olej. Dodatkowo obowiązuje techniczny limit 300 kroków (obcięcie epizodu),
który zabezpiecza przed nieskończonymi epizodami.

\subsection{Pojedynczy krok interakcji}
Agent obserwuje stan $s_t$ i wybiera akcję $a_t$. Środowisko wykonuje ruch,
wydobycie lub podparcie, zmniejsza zapas oleju, losuje zawał i sprawdza warunki
końca. Następnie zwraca nowy stan $s_{t+1}$, nagrodę $r_{t+1}$ oraz informację
o zakończeniu epizodu.

\subsection{Miary sukcesu}
Średni zwrot z epizodu, odsetek bezpiecznych powrotów, średnia wartość
wyniesionego ładunku, odsetek zawałów oraz średnia długość epizodu. Wyniki będą
porównywane z agentem losowym i z heurystyką progową („kop do osiągnięcia progu
ładunku, potem wracaj najkrótszą drogą”).

\subsection{Elementy do weryfikacji po pierwszych testach}
Wartości $p_0$, kary za zawał, zapasu oleju i liczby stempli, rozmiar map oraz
liczba złóż. Jeżeli przestrzeń stanów okaże się zbyt duża dla metod tablicowych,
rozważone zostanie zmniejszenie map lub zastosowanie DQN.

\section{Środowisko}
\doUzupelnienia{opis implementacji i interfejsu inicjalizacji oraz kroku,
legalnych akcji, map lub scenariuszy i zdarzeń dynamicznych.
Udokumentować rzeczywiście wykonane testy przypadków brzegowych.}

\section{Agent bazowy}
\doUzupelnienia{adekwatny agent losowy, heurystyczny lub inny baseline,
jego założenia i wyniki. Porównanie musi wykorzystywać uczciwe warunki.}

\section{Algorytm uczący się}
Porównywane są dwa tablicowe algorytmy uczenia różnicami czasowymi (TD):
Q-learning i SARSA \cite{sutton}. Oba utrzymują tablicę wartości akcji
$Q(s,a)$, przechowywaną jako słownik: kluczem jest stan, wartością wektor
$Q(s,\cdot)$ dla pięciu akcji. Pamięć zajmują więc tylko stany faktycznie
odwiedzone.

\textbf{Q-learning} (off-policy) aktualizuje wartość w kierunku najlepszej akcji
w stanie następnym:
\[
  Q(s,a) \leftarrow Q(s,a) + \alpha\left[r + \gamma \max_{a'} Q(s',a') - Q(s,a)\right].
\]
\textbf{SARSA} (on-policy) używa akcji $a'$ faktycznie wybranej przez bieżącą
politykę $\varepsilon$-zachłanną:
\[
  Q(s,a) \leftarrow Q(s,a) + \alpha\left[r + \gamma\, Q(s',a') - Q(s,a)\right].
\]
Ponieważ SARSA uwzględnia w ocenie także losowe ruchy eksploracyjne, w
środowiskach z ryzykiem uczy się zwykle ostrożniejszych strategii. W kopalni
oznacza to oczekiwanie, że SARSA będzie rzadziej schodzić głęboko przy niskiej
stabilności. Właśnie to porównanie motywuje wybór tych dwóch algorytmów.

\textbf{Szczegóły implementacji.} Akcje wybierane są polityką
$\varepsilon$-zachłanną z liniowo malejącym $\varepsilon$, a remisy wartości
rozstrzygane losowo. Stan terminalny (wyjście, zawał, brak oleju) nie jest
bootstrapowany. Obcięcie epizodu limitem kroków jest bootstrapowane, bo nie jest
prawdziwym końcem zadania. Faza uczenia jest oddzielona od ewaluacji: ewaluacja
używa polityki zachłannej ($\varepsilon=0$), bez aktualizacji tablicy Q i na
ziarnach losowych rozłącznych z treningiem. Implementacja:
\texttt{miner\_rl/agents/tabular.py}.

\textbf{Weryfikacja implementacji.} Przed użyciem w kopalni oba algorytmy
przetestowano na klasycznym zadaniu \emph{Cliff Walking} \cite{sutton}
(\texttt{tests/test\_tabular.py}). Zgodnie z literaturą Q-learning nauczył się
najkrótszej drogi wzdłuż urwiska (13 kroków), a SARSA dłuższej, bezpieczniejszej
(15--17 kroków). Testy sprawdzają także powtarzalność treningu przy ustalonym
ziarnie, brak bootstrapu w stanie terminalnym oraz zapis i odczyt tablicy Q.

\section{Plan eksperymentów}
\subsection{Porównywani agenci}
\begin{itemize}
  \item \textbf{Agent losowy} -- dolne odniesienie.
  \item \textbf{Heurystyka progowa} -- wydobywa najbliższe złoża, dopóki wartość
  ładunku nie osiągnie progu $X$, potem wraca najkrótszą drogą; stawia stempel,
  gdy stabilność spadnie do 1. Próg $X \in \{3, 6, 10, 15\}$ dobierany jest
  osobno dla każdej mapy na epizodach treningowych, a raportowany jest najlepszy
  wariant, aby porównanie było uczciwe.
  \item \textbf{Q-learning} i \textbf{SARSA} z tymi samymi hiperparametrami.
\end{itemize}

\subsection{Mapy}
Cztery mapy o rosnącej trudności: \emph{stabilna} (płytkie złoża, mała
głębokość), \emph{krucha} (szybki spadek bezpieczeństwa), \emph{kusząca głębia}
(najcenniejsze złoża na dnie) oraz \emph{zawał korytarza} (zdarzenie odcinające
najkrótszą drogę powrotu).

\subsection{Hiperparametry (wartości wyjściowe)}
\begin{center}
\begin{tabular}{@{}ll@{}}
\toprule
Parametr & Wartość \\
\midrule
współczynnik uczenia $\alpha$ & $0{,}1$ \\
dyskonto $\gamma$ & $0{,}99$ \\
$\varepsilon$ & liniowo od $1{,}0$ do $0{,}05$ w ciągu 70\% epizodów \\
liczba epizodów treningu & $20\,000$ na uruchomienie (do korekty po pierwszych testach) \\
liczba niezależnych uruchomień & 10 ziaren na każdą kombinację agent--mapa \\
\bottomrule
\end{tabular}
\end{center}

\subsection{Procedura i miary}
Co 500 epizodów treningu wykonywana jest ewaluacja zachłanna na 100 epizodach,
a po treningu ewaluacja końcowa na 1000 epizodach. Mierzone są: średni zwrot,
odsetek bezpiecznych powrotów, średnia wartość wyniesionego ładunku, odsetek
zawałów i wyczerpania oleju, średnia długość epizodu oraz największa osiągnięta
głębokość. Wyniki podawane są jako średnia i odchylenie standardowe z 10 ziaren.
Wszystkie ziarna, hiperparametry i konfiguracja środowiska zapisywane są razem
z wynikami (\texttt{results/}), co pozwala odtworzyć każdy wykres.

\subsection{Eksperymenty}
\begin{enumerate}
  \item \textbf{Porównanie główne.} Agent losowy, heurystyka progowa, Q-learning
  i SARSA na czterech mapach.
  \item \textbf{Wpływ kary za zawał.} Q-learning i SARSA z karą $-5$, $-20$ i
  $-50$ (10 ziaren każda) na mapie \emph{kusząca głębia}. Badana jest zmiana
  skłonności do ryzyka: głębokość wydobycia, odsetek zawałów i wartość ładunku.
\end{enumerate}

\subsection{Hipotezy}
\begin{itemize}
  \item[H1.] Agenci RL osiągają wyższy średni zwrot niż heurystyka progowa,
  zwłaszcza na mapach, gdzie opłaca się zmienić kolejność wydobycia lub użyć stempla.
  \item[H2.] SARSA w trakcie treningu powoduje mniej zawałów niż Q-learning,
  kosztem niższej wartości ładunku. W ewaluacji zachłannej różnica może maleć.
  \item[H3.] Wyższa kara za zawał prowadzi do płytszego kopania i rzadszych
  zawałów. Zbyt niska kara skutkuje zbyt ryzykowną strategią.
\end{itemize}

\subsection{Planowane wykresy i tabele}
Krzywe uczenia (średnia z pasmem odchylenia z 10 ziaren, oś X: epizod, oś Y:
zwrot) oraz krzywe odsetka zawałów w trakcie treningu. Tabela wyników końcowych
dla wszystkich agentów i map. Wykres zależności odsetka zawałów i wartości
ładunku od kary. Mapa cieplna odwiedzin pól dla wybranych polityk. Analiza
nieudanych epizodów: zestawienie przyczyn porażek oraz nagrania wybranych
epizodów.

\section{Wyniki}
Nie przeprowadzono jeszcze eksperymentów. Nie ma wyników liczbowych do przedstawienia.
\doUzupelnienia{wyniki odtwarzalne z kodu lub danych, przebieg uczenia,
porównania, zmienność wyników i przypadki niepowodzeń.
Wykresy muszą mieć opisane osie, jednostki i legendy.}

\section{Analiza i wnioski}
\doUzupelnienia{interpretacja faktycznych wyników, odpowiedzi na hipotezy,
co poprawia lub pogarsza wyniki oraz własne wnioski studenta.}

\section{Ograniczenia rozwiązania}
Na obecnym etapie ograniczeniem jest brak implementacji i wyników.
\doUzupelnienia{ograniczenia modelu, obserwacji, funkcji nagrody, eksperymentów
i możliwości uogólniania, zgodne z rzeczywistym stanem rozwiązania.}

\section{Instrukcja uruchomienia}
Repozytorium:
\url{https://github.com/Kesacky/uczenie-ze-wzmocnieniem}.

Nie ma jeszcze programu RL do uruchomienia.
Dokumentację można kompilować poleceniem \texttt{pdflatex main.tex}
z katalogu \texttt{report}; kompilację należy wykonać dwukrotnie,
aby uaktualnić spis treści.
\doUzupelnienia{wymagania środowiska programu, instalacja zależności,
polecenia uruchomienia, treningu, ewaluacji i odtworzenia wyników
z czystej kopii repozytorium.}

\section{Wsparcie AI}
\subsection*{21 października 2026 - przygotowanie projektu}
\textbf{Narzędzie i model:} Codex, rodzina GPT-6; dokładny wariant modelu
nie został udostępniony w interfejsie tej sesji. Korzystano także z podagentów
w ramach Codex do analizy wymagań i przygotowania narzędzia dziennika.

\textbf{Cel:} odczytanie dwóch dokumentów kursowych i przygotowanie
repozytorium oraz dokumentacji wymaganej do dalszej realizacji projektu.

\textbf{Zakres pomocy:} analiza wymagań P00 i P01, zestawienie terminów
i list kontrolnych, przygotowanie struktury repozytorium i szablonu raportu
LaTeX, narzędzia zapisu promptów i odpowiedzi oraz sprawdzenie dokumentacji.
Na tym etapie AI nie opracowało algorytmu RL, środowiska ani wyników eksperymentów.

\textbf{Weryfikacja studenta:} niepotwierdzona. Testy i kontrole wykonane
przez agenta nie stanowią oświadczenia o samodzielnej weryfikacji studenta.
Po zapoznaniu się z materiałami student powinien wskazać, które elementy
sam sprawdził, zrozumiał lub zmodyfikował.

\textbf{Transkrypt:} pełne widoczne prompty i odpowiedzi są zapisywane lokalnie
w dzienniku \texttt{private/ai-journal/main.tex} wraz z tekstami źródłowymi
i metadanymi. Sposób udostępnienia transkryptów prowadzącemu wymaga ustalenia.
Ta rozmowa odbywa się w narzędziu programistycznym; wymaganie dodatkowego
linku dla rozmów poza takim narzędziem nie ma tu zastosowania.

\textbf{Dalsze aktualizacje:} każdą kolejną rozmowę dotyczącą projektu należy
zapisać w całości oraz uzupełnić niniejszą sekcję o datę, model, cel, zakres
wykorzystania i rzeczywisty status samodzielnej weryfikacji studenta.


\begin{thebibliography}{9}
\bibitem{sutton}
R.~S. Sutton, A.~G. Barto,
\emph{Reinforcement Learning: An Introduction}, wyd.~2, MIT Press, 2018.
\bibitem{zasady}
Materiały kursowe WEiI PRz,
\emph{Uczenie się ze wzmocnieniem. Zasady realizacji projektu},
rok akademicki 2026/2027, w szczególności sekcje 2, 4, 8 i 9.
Materiał źródłowy przechowywany lokalnie.
\bibitem{checklista}
Materiały kursowe WEiI PRz,
\emph{Uczenie się ze wzmocnieniem. Lista kontrolna projektu},
rok akademicki 2026/2027, w szczególności sekcje 2--6.
Materiał źródłowy przechowywany lokalnie.
\end{thebibliography}
\end{document}
